% Working technical note for the compositional AND-gate case study.
% Kept separate from the paper source for model/theory review.

\section{Compositional AND-gate models and contract construction}
\label{sec:and_gate_model_notes}

\subsection{Top-level specification and monotonicity}

Let
\[
L=[0.0055,0.07],\qquad H=[0.35,0.65]
\]
denote LOW and HIGH promoter-activity intervals in RPU. The current top-level
specification is
\begin{align}
(\xi_1,\xi_2)\in H\times H &\Longrightarrow y\ge 120~{\rm nM/min},\\
(\xi_1,\xi_2)\in L\times H &\Longrightarrow y\le 60~{\rm nM/min},\\
(\xi_1,\xi_2)\in H\times L &\Longrightarrow y\le 60~{\rm nM/min}.
\end{align}
A separate LOW--LOW mode is unnecessary once coordinatewise monotonicity is
established. If
\[
\xi_1\le\xi_1',\quad \xi_2\le\xi_2'
\Longrightarrow
y(\xi_1,\xi_2)\le y(\xi_1',\xi_2'),
\]
then for any $(\xi_1,\xi_2)\in L\times L$ and any
$\widehat\xi_2\in H$,
\[
y(\xi_1,\xi_2)\le y(\xi_1,\widehat\xi_2)\le 60.
\]
Thus the LOW--LOW truth-table condition follows from the LOW--HIGH OFF
contract. The same argument follows from HIGH--LOW. The transition interval
$0.07<\xi_i<0.35$ remains unspecified.

\subsection{Lumped expression modules and the common complex model}

The implementation uses
\[
Y_{i,\rm tot}=g_i\xi_i.
\]
A minimal ODE realization is
\[
\dot Y_{i,\rm tot}=\alpha_i\xi_i-\delta_iY_{i,\rm tot},
\qquad
g_i=\frac{\alpha_i}{\delta_i},
\]
whose constant-input steady state gives the implemented gain relation.

For the initial designs, the two expressed partners form an active complex:
\begin{align}
\dot C&=k_f(Y_{1,\rm tot}-C)(Y_{2,\rm tot}-C)-k_rC,\\
\dot P&=k_{\rm out}C,\qquad y:=\dot P.
\end{align}
At quasi-steady binding,
\[
KC=(Y_{1,\rm tot}-C)(Y_{2,\rm tot}-C),
\qquad K=\frac{k_r}{k_f},
\]
and
\[
\frac{Y_{1,\rm tot}Y_{2,\rm tot}}{K+Y_{1,\rm tot}+Y_{2,\rm tot}}
\le C\le
\min\left\{Y_{1,\rm tot},Y_{2,\rm tot},
\frac{Y_{1,\rm tot}Y_{2,\rm tot}}{K}\right\}.
\]
The direct core is increasing in both component abundances:
\[
\frac{\partial C}{\partial Y_1}
=
\frac{Y_2-C}{K+Y_1+Y_2-2C}>0,\qquad
\frac{\partial C}{\partial Y_2}
=
\frac{Y_1-C}{K+Y_1+Y_2-2C}>0,
\]
and decreasing in $K$:
\[
\frac{\partial C}{\partial K}
=
-\frac{C}{K+Y_1+Y_2-2C}<0.
\]
These properties justify the corner evaluations used in the current
finite-region calculations.

\subsection{Design 1: HrpR/HrpS with tunable HrpV inhibition}

The initial effective parameters are
\[
K=250~{\rm nM},\quad k_{\rm out}=1.2~{\rm min}^{-1},\quad
g_1=4000~{\rm nM/RPU},\quad g_2=800~{\rm nM/RPU}.
\]
The redesigned computation uses
\[
C_{\rm act}=[C-V_{\rm eff}]_+,\qquad
y=k_{\rm out}C_{\rm act},
\]
where $V_{\rm eff}$ is the tunable HrpV regulation level.

A compatible ODE realization introduces an inhibited complex $I$:
\begin{align}
\dot Q&=k_f(R_{\rm tot}-Q)(S_{\rm tot}-Q)-k_rQ,\\
\dot I&=k_v^+(Q-I)(V_{\rm tot}-I)-k_v^-I,\\
\dot P&=k^ý\rm out}(Q-I).
\end{align}
In the fast, strong-inhibition limit,
\[
Q-I\approx[Q-V_{\rm tot}]_+,
\]
which recovers the effective model implemented in the code. Since the direct
complex $C$ is increasing in both inputs and the positive-part map is
nondecreasing, the HrpV redesign is coordinatewise nondecreasing.

Contract satisfaction gives
\[
V_{\rm eff}\ge
\max\left\{
C_{\rm OFF,1}-\frac{y_{\rm off}}{k^ý\rm out}},
C_{\rm OFF,2}-\frac{y_{\rm off}}{k_{\rm out}},0
\right\},
\]
and
\[
V_{\rm eff}\le
C_{\rm ON}-\frac{y_{\rm on}}{k^ý\rm out}}.
\]
The current computation gives approximately
\[
114.4~{\rm nM}\le V_{\rm eff}\le130.7~{\rm nM},
\]
and uses $V_{\rm eff}=120~{\rm nM}$.

\subsection{Design 2: split T7 RNA polymerase}

The initial effective parameters are
\[
K=1200~{\rm nM},\quad k_{\rm out}=10~{\rm min}^{-1},\quad
g_1=500~{\rm nM/RPU},\quad g_2=700~{\rm nM/RPU}.
\]
The redesign retains the same ODE/QSS structure and uses the effective
fragment-association parameter $K$ as the design variable. For a required
complex level $C^\star$,
\[
K=\frac{(Y_1-C^\star)(Y_2-C^\star)}{C^\star}.
\]
The ON requirement gives an upper bound on $K$ and the OFF requirements give
lower bounds. The current fixed-parameter calculation gives
\[
2.286~\mu{\rm M}\lesssim K\lesssim3.165~\mu{\rm M},
\]
and uses $K=3.0~\mu{\rm M}$.

\subsection{Design 3: InvF/SicA with SicA--SipC turnover}

The initial effective parameters are
\[
K^ýXY}=250~{\rm nM},\quad k_{\rm out}=1~{\rm min}^{-1},\quad
g_1=800~{\rm nM/RPU},\quad g_2=4000~{\rm nM/RPU}.
\]
The initial design uses the direct complex model above.

For the redesigned model, let $X$ be free InvF, $Y$ free SicA, $Z$ free
SipC, $C_{XY}$ the active InvF:SicA complex, and $C_{YZ}$ the SicA:SipC
complex. A full ODE realization consistent with the reduced steady-state
calculation is
\begin{align}
\dot X&=p_X-\gamma_XX-k_{XY}^+XY+k^ýXY}^-C_{XY},\\
\dot Y&=p_Y-\gamma_YY-k_{XY}^+XY+k^ýXY}^-C_{XY}
-k_{YZ}^+YZ+(k_{YZ}^-+k_Z)C_{YZ},\\
\dot Z&=p_Z-\gamma_ZZ-k_{YZ}^+YZ+k^ýYZ}^-C_{YZ},\\
\dot C_{XY}&=k_{XY}^+XY-k^ýXY}^-C_{XY}-\gamma_{XY}C_{XY},\\
\dot C_{YZ}&=k_{YZ}^+YZ-k^ýYZ}^-C_{YZ}
-(\gamma_{YZ}+k_Z)C_{YZ},\\
\dot P&=k^ý\rm out}C_{XY}.
\end{align}
The $k_ZC_{YZ}$ term releases SicA while removing the SipC partner.

The implemented computation uses the fast-binding relations
\[
C_{XY}=\frac{XY}{K^ýXY}},\qquad
C_{YZ}=\frac{YZ}{K_{YZ}},
\]
and steady-state balances
\begin{align}
p_X&=\gamma_XX+\gamma_{XY}C_{XY},\\
p_Y&=\gamma_YY+\gamma_{XY}C_{XY}+\gamma_{YZ}C_{YZ},\\
p_Z&=\gamma_ZZ+(\gamma_{YZ}+k_Z)C_{YZ}.
\end{align}
The input-dependent production terms are
\[
p_X=\gamma_Xg_1\xi_1,\qquad
p_Y=\gamma_Yg_2\xi_2.
\]
The parameters used in the current code are
\[
K_{YZ}=28~{\rm nM},\quad
\gamma_X=\gamma_{XY}=0.022~{\rm min}^{-1},\quad
\gamma_Y=0.046~{\rm min}^{-1},
\]
\[
\gamma_Z=0.139~{\rm min}^{-1},\quad
\gamma_{YZ}=0.022~{\rm min}^{-1},\quad
k_Z=0.1~{\rm min}^{-1},
\]
together with the existing effective $K_{XY}=250~{\rm nM}$ and
$k_{\rm out}=1~{\rm min}^{-1}$.

\paragraph{Monotonicity of the reduced InvF/SicA model.}
Define
\[
h(Y)=\frac{Y}{K^ýXY}\gamma_X+\gamma_{XY}Y},\qquad
q(Y)=\frac{Y}{K_{YZ}\gamma_Z+(\gamma_{YZ}+k_Z)Y}.
\]
Then
\[
C_{XY}=p_Xh(Y),\qquad C_{YZ}=p_Zq(Y),
\]
and the SicA balance is
\[
R(Y;p_X,p_Y)
=
\gamma_YY+\gamma_{XY}p_Xh(Y)
+\gamma_{YZ}p_Zq(Y)-p_Y=0.
\]
Since $h'(Y)>0$ and $q'(Y)>0$,
\[
R_Y=
\gamma_Y+\gamma_{XY}p_Xh'(Y)+\gamma_{YZ}p_Zq'(Y)>0.
\]
Hence the reduced steady state is unique and
\[
\frac{\partial Y^*}{\partial p_Y}=\frac{1}{R_Y}>0,\qquad
\frac{\partial Y^*}{\partial p_X}
=-\frac{\gamma_{XY}h(Y^*)}{R_Y}<0.
\]
The active complex satisfies
\[
\frac{\partial C_{XY}^*}{\partial p_Y}
=
\frac{p_Xh'(Y^*)}{R_Y}>0,
\]
and
\[
\frac{\partial C_{XY}^*}{\partial p_X}
=
h(Y^*)
\frac{\gamma_Y+\gamma_{YZ}p_Zq'(Y^*)}{R_Y}>0.
\]
Thus the reduced redesigned output is coordinatewise nondecreasing in
$(\xi_1,\xi_2)$, which justifies its corner-based mode certification.

\subsection{Finite-box contract generation}

The current code constructs a local core contract for each top-level input
region. This can be formalized as a contract abstraction operator
\[
\mathfrak C^ý\Sigma}(B,\Theta),
\]
where $\Sigma$ is the subsystem, $B=[\underline\xi,\overline\xi]$ is a finite
input box, and $\Theta$ is the parameter set used for verification.

For a coordinatewise monotone input-output map,
\[
y_{\min}(B,\theta)=y(\underline\xi,\theta),\qquad
y_{\max}(B,\theta)=y(\overline\xi,\theta).
\]
If parameter monotonicity is known, adverse parameter endpoints can be selected
in the same way:
\[
\underline y(B,\Theta)
=
\inf_{\theta\in\Theta}y(\underline\xi,\theta),\qquad
\overline y(B,\Theta)
=
\sup_{\theta\in\Theta}y(\overline\xi,\theta).
\]
These extrema give a sound local contract over $B$. For the direct QSS core,
the structural inequalities on $C$ provide additional affine log-coordinate
guarantees.

The implementation therefore generates
\[
\mathfrak C_{\Sigma_3}(B_{\rm ON},\Theta),\quad
\mathfrak C_{\Sigma_3}(B_{\rm OFF,1},\Theta),\quad
\mathfrak C_{\Sigma_3}(B_{\rm OFF,2},\Theta).
\]
This is one biochemical subsystem with three local abstractions. An alternative
would be one core contract over a larger box containing all three environments;
that contract would generally be looser.

\subsection{Fixed-parameter redesign verification}

The initial designs use provisional multiplicative parameter ranges
\[
f_g=1.25,\qquad f_K=1.30,\qquad f_r=1.25.
\]
The redesigned calculations currently set those factors to one and evaluate
the selected redesign parameters at fixed values.

The interpretation is: for each redesigned architecture, the contract is
verified for every input in the specified ON or OFF input region at the
selected parameter values. The input contract still covers a continuum of
promoter activities. Only the parameter set is a single selected design point.
This is sufficient for the current contract-guided parameter-synthesis
question. Parameter robustness is a separate question and is not required for
the present claim.

The phrase ``singleton contract'' should therefore be avoided. ``Fixed-parameter
verification'' or ``verification at the selected synthesis point'' describes
the calculation more accurately.

\subsection{Zero output and the logarithmic epsilon floor}

The HrpV effective model can produce exactly zero. Since the implementation
uses $z_y=\log y$, zero has no finite log-coordinate representation. The code
therefore uses
\[
\varepsilon_y=10^{-12}~{\rm nM/min}
\]
only when a zero upper bound must be represented in log coordinates. Thus an
exact physical-model statement $y=0$ is encoded computationally as
\[
y\le\varepsilon_y.
\]
This is a numerical representation floor. It is not a biological leak model
and is not an additional physical parameter.

\subsection{Numerical assumption padding}

The current implementation expands the local core assumption interval by the
factor $1.001$. The code comment states that this was introduced to avoid
numerical equality artifacts when Pacti eliminates internal
component-abundance variables. In exact contract arithmetic, equality of an
upstream guarantee boundary and a downstream assumption boundary is admissible.

An outward pad is formally problematic if the core output bounds are derived
on the original interval and then asserted over the larger padded interval.
The preferred implementation is therefore either
\begin{enumerate}
\item use the exact physical assumption bounds if Pacti handles the shared
boundaries reliably; or
\item if a numerical margin is necessary, define it explicitly as a numerical
tolerance and recompute the core guarantees on the enlarged interval.
\end{enumerate}
The need for the present $1.001$ factor should be tested before it is retained.

\subsection{Theory development}

The application suggests four closely connected formal results.

\paragraph{Monotone CRN-to-contract abstraction.}
For a CRN subsystem whose input-output map is monotone on a finite orthant box,
derive a sound polyhedral contract from order extrema:
\[
(\Sigma,B,\Theta)\longmapsto \mathfrak C_\Sigma(B,\Theta).
\]

\paragraph{Monotonicity closure of the AND specification.}
For coordinatewise nondecreasing two-input implementations, the LOW--HIGH and
HIGH--LOW OFF contracts imply LOW--LOW OFF behavior. This gives a formal reason
for the three-mode top-level specification used here.

\paragraph{Contract-guided parameter synthesis.}
For a design parameter $p$, define
\[
\mathcal P_{\rm feasible}
=
\left\{
p:
\mathfrak C_{\Sigma(p)}(B_q,\Theta_p)
\preceq C_{\rm spec}^q
\ \text{for every mode }q
\right\}.
\]
The HrpV interval and split-T7 $K$ interval are explicit examples in which this
set is obtained analytically from monotonicity and the top-level thresholds.

\paragraph{Mechanism-specific contract templates.}
The direct binding core, positive-part inhibited core, and InvF/SicA turnover
core admit different analytic bounding structures. A general result can
identify CRN motifs for which finite-box contracts can be generated
automatically and remain polyhedral after the logarithmic transformation.

The monotonicity arguments in Incer, Pandey, Nolan, Peterman, Galloway, Sontag,
and Del Vecchio (CDC 2024) provide the natural prior formal starting point.
The next theory pass should distinguish the monotonicity machinery inherited
from that work from the new finite-box CRN-to-contract construction,
monotonicity closure result, and refinement-based parameter-synthesis results
developed here. The exact theorem-level comparison should be added after
reviewing that specific paper in detail.
